\appendix
\chapter{Mathematical Proofs and Derivations}

This appendix provides rigorous mathematical foundations for the statistical methods used in A/B testing. We present detailed proofs and derivations with step-by-step explanations.

\section{Sample Size Formulas}

Sample size determination is fundamental to experimental design. We derive formulas for various testing scenarios.

\subsection{Sample Size for Two-Sample t-Test}

\begin{theorem}[Sample Size for Comparing Two Means]
To detect a difference $\delta = \mu_2 - \mu_1$ between two population means with common variance $\sigma^2$, significance level $\alpha$, and power $1-\beta$, the required sample size per group is:
\begin{equation}
n = \frac{2(z_{\alpha/2} + z_\beta)^2 \sigma^2}{\delta^2}
\end{equation}
where $z_{\alpha/2}$ and $z_\beta$ are the critical values from the standard normal distribution.
\end{theorem}

\begin{proof}
Consider testing $H_0: \mu_1 = \mu_2$ versus $H_1: \mu_1 \neq \mu_2$ using the test statistic:
\begin{equation}
T = \frac{\bar{X}_2 - \bar{X}_1}{\sqrt{2\sigma^2/n}}
\end{equation}

\textbf{Step 1: Distribution under the null hypothesis.}

Under $H_0$, the difference $\bar{X}_2 - \bar{X}_1$ has mean 0 and variance:
\begin{equation}
\text{Var}(\bar{X}_2 - \bar{X}_1) = \text{Var}(\bar{X}_2) + \text{Var}(\bar{X}_1) = \frac{\sigma^2}{n} + \frac{\sigma^2}{n} = \frac{2\sigma^2}{n}
\end{equation}

By the Central Limit Theorem, for large $n$:
\begin{equation}
T \mid H_0 \sim \mathcal{N}(0, 1)
\end{equation}

\textbf{Step 2: Rejection region.}

For a two-sided test at level $\alpha$, we reject $H_0$ if $|T| > z_{\alpha/2}$, which occurs when:
\begin{equation}
|\bar{X}_2 - \bar{X}_1| > z_{\alpha/2} \sqrt{\frac{2\sigma^2}{n}}
\end{equation}

\textbf{Step 3: Distribution under the alternative.}

Under $H_1: \mu_2 - \mu_1 = \delta$, the test statistic has mean:
\begin{align}
\mathbb{E}[T \mid H_1] &= \frac{\mathbb{E}[\bar{X}_2 - \bar{X}_1]}{\sqrt{2\sigma^2/n}} \\
&= \frac{\delta}{\sqrt{2\sigma^2/n}} \\
&= \frac{\delta\sqrt{n}}{\sqrt{2}\sigma}
\end{align}

and variance 1. Therefore:
\begin{equation}
T \mid H_1 \sim \mathcal{N}\left(\frac{\delta\sqrt{n}}{\sqrt{2}\sigma}, 1\right)
\end{equation}

\textbf{Step 4: Power calculation.}

The power is the probability of rejecting $H_0$ when $H_1$ is true. Assuming $\delta > 0$:
\begin{align}
\text{Power} &= P(T > z_{\alpha/2} \mid H_1) \\
&= P\left(\mathcal{N}\left(\frac{\delta\sqrt{n}}{\sqrt{2}\sigma}, 1\right) > z_{\alpha/2}\right) \\
&= P\left(\mathcal{N}(0,1) > z_{\alpha/2} - \frac{\delta\sqrt{n}}{\sqrt{2}\sigma}\right) \\
&= \Phi\left(\frac{\delta\sqrt{n}}{\sqrt{2}\sigma} - z_{\alpha/2}\right)
\end{align}

where $\Phi$ is the standard normal CDF.

\textbf{Step 5: Solve for n.}

For power $1-\beta$, we require:
\begin{equation}
\Phi\left(\frac{\delta\sqrt{n}}{\sqrt{2}\sigma} - z_{\alpha/2}\right) = 1 - \beta
\end{equation}

Since $\Phi(z_\beta) = 1 - \beta$, we have:
\begin{equation}
\frac{\delta\sqrt{n}}{\sqrt{2}\sigma} - z_{\alpha/2} = z_\beta
\end{equation}

Solving for $n$:
\begin{align}
\frac{\delta\sqrt{n}}{\sqrt{2}\sigma} &= z_{\alpha/2} + z_\beta \\
\sqrt{n} &= \frac{\sqrt{2}\sigma(z_{\alpha/2} + z_\beta)}{\delta} \\
n &= \frac{2(z_{\alpha/2} + z_\beta)^2 \sigma^2}{\delta^2}
\end{align}
\end{proof}

\subsection{Sample Size for Comparing Two Proportions}

\begin{theorem}[Sample Size for Proportion Comparison]
To detect a difference between two proportions $p_1$ and $p_2$ with significance level $\alpha$ and power $1-\beta$, the required sample size per group is:
\begin{equation}
n = \frac{(z_{\alpha/2} + z_\beta)^2 [p_1(1-p_1) + p_2(1-p_2)]}{(p_2 - p_1)^2}
\end{equation}
\end{theorem}

\begin{proof}
\textbf{Step 1: Test statistic under the alternative.}

Under $H_1: p_1 \neq p_2$, the difference in sample proportions has:
\begin{align}
\mathbb{E}[\hat{p}_2 - \hat{p}_1] &= p_2 - p_1 \\
\text{Var}(\hat{p}_2 - \hat{p}_1) &= \frac{p_2(1-p_2)}{n} + \frac{p_1(1-p_1)}{n} = \frac{p_1(1-p_1) + p_2(1-p_2)}{n}
\end{align}

Define the standardized statistic:
\begin{equation}
Z = \frac{(\hat{p}_2 - \hat{p}_1) - (p_2 - p_1)}{\sqrt{[p_1(1-p_1) + p_2(1-p_2)]/n}} \sim \mathcal{N}(0, 1)
\end{equation}

\textbf{Step 2: Test statistic under the null.}

Under $H_0: p_1 = p_2 = p$ (pooled proportion), the variance is:
\begin{equation}
\text{Var}_0(\hat{p}_2 - \hat{p}_1) = \frac{2p(1-p)}{n}
\end{equation}

The null statistic is:
\begin{equation}
Z_0 = \frac{\hat{p}_2 - \hat{p}_1}{\sqrt{2p(1-p)/n}} \sim \mathcal{N}(0, 1)
\end{equation}

\textbf{Step 3: Critical value.}

Reject $H_0$ if $|Z_0| > z_{\alpha/2}$, which occurs when:
\begin{equation}
|\hat{p}_2 - \hat{p}_1| > z_{\alpha/2}\sqrt{\frac{2p(1-p)}{n}}
\end{equation}

\textbf{Step 4: Power under the alternative.}

Assuming $p_2 > p_1$ and approximating $p \approx \frac{p_1 + p_2}{2}$, the power is:
\begin{align}
\text{Power} &= P\left(\hat{p}_2 - \hat{p}_1 > z_{\alpha/2}\sqrt{\frac{2p(1-p)}{n}} \mid H_1\right) \\
&= P\left(\frac{(\hat{p}_2 - \hat{p}_1) - (p_2-p_1)}{\sqrt{[p_1(1-p_1) + p_2(1-p_2)]/n}} > \frac{z_{\alpha/2}\sqrt{2p(1-p)/n} - (p_2-p_1)}{\sqrt{[p_1(1-p_1) + p_2(1-p_2)]/n}}\right)
\end{align}

For large $n$, $2p(1-p) \approx p_1(1-p_1) + p_2(1-p_2)$, so:
\begin{equation}
\text{Power} \approx \Phi\left(\frac{(p_2-p_1)\sqrt{n}}{\sqrt{p_1(1-p_1) + p_2(1-p_2)}} - z_{\alpha/2}\right)
\end{equation}

\textbf{Step 5: Solve for n.}

Setting Power $= 1 - \beta$:
\begin{align}
\frac{(p_2-p_1)\sqrt{n}}{\sqrt{p_1(1-p_1) + p_2(1-p_2)}} - z_{\alpha/2} &= z_\beta \\
\sqrt{n} &= \frac{(z_{\alpha/2} + z_\beta)\sqrt{p_1(1-p_1) + p_2(1-p_2)}}{p_2 - p_1} \\
n &= \frac{(z_{\alpha/2} + z_\beta)^2[p_1(1-p_1) + p_2(1-p_2)]}{(p_2 - p_1)^2}
\end{align}
\end{proof}

\subsection{Minimum Detectable Effect}

\begin{corollary}[Minimum Detectable Effect]
Given fixed sample size $n$, significance level $\alpha$, and power $1-\beta$, the minimum detectable effect for comparing means is:
\begin{equation}
\text{MDE} = \delta = \frac{(z_{\alpha/2} + z_\beta)\sigma\sqrt{2}}{\sqrt{n}}
\end{equation}
\end{corollary}

\begin{proof}
Rearranging the sample size formula:
\begin{align}
n &= \frac{2(z_{\alpha/2} + z_\beta)^2 \sigma^2}{\delta^2} \\
\delta^2 &= \frac{2(z_{\alpha/2} + z_\beta)^2 \sigma^2}{n} \\
\delta &= \frac{(z_{\alpha/2} + z_\beta)\sigma\sqrt{2}}{\sqrt{n}}
\end{align}
\end{proof}

\subsection{Unequal Sample Sizes}

\begin{theorem}[Sample Size with Allocation Ratio]
If the allocation ratio is $r = n_2/n_1$ (i.e., $n_2 = rn_1$), the required sample size for group 1 is:
\begin{equation}
n_1 = \frac{(z_{\alpha/2} + z_\beta)^2 \sigma^2(1 + 1/r)}{\delta^2}
\end{equation}
and $n_2 = rn_1$.
\end{theorem}

\begin{proof}
The variance of the difference with unequal sample sizes is:
\begin{equation}
\text{Var}(\bar{X}_2 - \bar{X}_1) = \frac{\sigma^2}{n_2} + \frac{\sigma^2}{n_1} = \frac{\sigma^2}{rn_1} + \frac{\sigma^2}{n_1} = \frac{\sigma^2(1 + r)}{rn_1}
\end{equation}

Following the same derivation as before, the test statistic under $H_1$ is:
\begin{equation}
T = \frac{(\bar{X}_2 - \bar{X}_1) - \delta}{\sqrt{\sigma^2(1+r)/(rn_1)}} \sim \mathcal{N}(0,1)
\end{equation}

The power condition gives:
\begin{equation}
\frac{\delta}{\sqrt{\sigma^2(1+r)/(rn_1)}} - z_{\alpha/2} = z_\beta
\end{equation}

Solving for $n_1$:
\begin{align}
\frac{\delta\sqrt{rn_1}}{\sigma\sqrt{1+r}} &= z_{\alpha/2} + z_\beta \\
n_1 &= \frac{(z_{\alpha/2} + z_\beta)^2\sigma^2(1+r)}{r\delta^2} = \frac{(z_{\alpha/2} + z_\beta)^2\sigma^2(1 + 1/r)}{\delta^2}
\end{align}
\end{proof}

\begin{remark}
The total sample size $N = n_1 + n_2 = n_1(1+r)$ is minimized when $r = 1$ (equal allocation).
\end{remark}

\section{Central Limit Theorem}

The Central Limit Theorem is the foundation of large-sample inference in A/B testing.

\subsection{Classical Central Limit Theorem}

\begin{theorem}[Lindeberg-Lévy CLT]
Let $X_1, X_2, \ldots$ be independent and identically distributed random variables with $\mathbb{E}[X_i] = \mu$ and $\text{Var}(X_i) = \sigma^2 < \infty$. Define $S_n = \sum_{i=1}^n X_i$ and $\bar{X}_n = S_n/n$. Then:
\begin{equation}
\frac{S_n - n\mu}{\sigma\sqrt{n}} = \frac{\bar{X}_n - \mu}{\sigma/\sqrt{n}} \xrightarrow{d} \mathcal{N}(0, 1)
\end{equation}
as $n \to \infty$, where $\xrightarrow{d}$ denotes convergence in distribution.
\end{theorem}

\begin{proof}[Proof Sketch via Characteristic Functions]
We prove this using moment generating functions (characteristic functions).

\textbf{Step 1: Standardize the sum.}

Define $Z_i = (X_i - \mu)/\sigma$ so that $\mathbb{E}[Z_i] = 0$ and $\text{Var}(Z_i) = 1$. Then:
\begin{equation}
\frac{S_n - n\mu}{\sigma\sqrt{n}} = \frac{1}{\sqrt{n}}\sum_{i=1}^n Z_i
\end{equation}

\textbf{Step 2: Characteristic function.}

Let $\varphi_Z(t) = \mathbb{E}[e^{itZ_1}]$ be the characteristic function of $Z_1$. The characteristic function of $\frac{1}{\sqrt{n}}\sum_{i=1}^n Z_i$ is:
\begin{equation}
\varphi_n(t) = \mathbb{E}\left[\exp\left(it\frac{1}{\sqrt{n}}\sum_{i=1}^n Z_i\right)\right] = \prod_{i=1}^n \mathbb{E}\left[e^{itZ_i/\sqrt{n}}\right] = \left[\varphi_Z\left(\frac{t}{\sqrt{n}}\right)\right]^n
\end{equation}

\textbf{Step 3: Taylor expansion.}

Since $\mathbb{E}[Z_1] = 0$ and $\text{Var}(Z_1) = 1$, we can expand:
\begin{align}
\varphi_Z(t) &= \mathbb{E}[e^{itZ_1}] \\
&= \mathbb{E}\left[1 + itZ_1 + \frac{(it)^2Z_1^2}{2!} + \frac{(it)^3Z_1^3}{3!} + \cdots\right] \\
&= 1 + it\mathbb{E}[Z_1] + \frac{(it)^2\mathbb{E}[Z_1^2]}{2} + o(t^2) \\
&= 1 - \frac{t^2}{2} + o(t^2)
\end{align}

\textbf{Step 4: Substitute and take limit.}

\begin{align}
\varphi_n(t) &= \left[\varphi_Z\left(\frac{t}{\sqrt{n}}\right)\right]^n \\
&= \left[1 - \frac{t^2}{2n} + o\left(\frac{1}{n}\right)\right]^n \\
&\to e^{-t^2/2} \quad \text{as } n \to \infty
\end{align}

by the standard limit $\lim_{n\to\infty}(1 + x/n)^n = e^x$.

\textbf{Step 5: Recognize the limit.}

The function $e^{-t^2/2}$ is the characteristic function of $\mathcal{N}(0,1)$.

By Lévy's continuity theorem, convergence of characteristic functions implies convergence in distribution:
\begin{equation}
\frac{S_n - n\mu}{\sigma\sqrt{n}} \xrightarrow{d} \mathcal{N}(0, 1)
\end{equation}
\end{proof}

\subsection{Berry-Esseen Theorem}

The Berry-Esseen theorem provides a quantitative bound on the rate of convergence in the CLT.

\begin{theorem}[Berry-Esseen Bound]
Under the conditions of the CLT, if additionally $\mathbb{E}[|X_i|^3] = \rho < \infty$, then:
\begin{equation}
\sup_{x \in \mathbb{R}} \left|P\left(\frac{\bar{X}_n - \mu}{\sigma/\sqrt{n}} \leq x\right) - \Phi(x)\right| \leq \frac{C\rho}{\sigma^3\sqrt{n}}
\end{equation}
where $\Phi$ is the standard normal CDF and $C$ is an absolute constant (best known value: $C \approx 0.4748$).
\end{theorem}

\begin{remark}
This shows convergence rate is $O(n^{-1/2})$. For practical purposes, $n \geq 30$ often suffices for near-normality.
\end{remark}

\subsection{Delta Method}

\begin{theorem}[Delta Method]
Let $\sqrt{n}(\hat{\theta}_n - \theta) \xrightarrow{d} \mathcal{N}(0, \sigma^2)$ and let $g: \mathbb{R} \to \mathbb{R}$ be differentiable at $\theta$ with $g'(\theta) \neq 0$. Then:
\begin{equation}
\sqrt{n}(g(\hat{\theta}_n) - g(\theta)) \xrightarrow{d} \mathcal{N}(0, [g'(\theta)]^2\sigma^2)
\end{equation}
\end{theorem}

\begin{proof}
By first-order Taylor expansion around $\theta$:
\begin{equation}
g(\hat{\theta}_n) = g(\theta) + g'(\theta)(\hat{\theta}_n - \theta) + R_n
\end{equation}
where the remainder satisfies $R_n = o_p(|\hat{\theta}_n - \theta|)$.

Rearranging:
\begin{equation}
g(\hat{\theta}_n) - g(\theta) = g'(\theta)(\hat{\theta}_n - \theta) + R_n
\end{equation}

Multiplying by $\sqrt{n}$:
\begin{equation}
\sqrt{n}(g(\hat{\theta}_n) - g(\theta)) = g'(\theta)\sqrt{n}(\hat{\theta}_n - \theta) + \sqrt{n}R_n
\end{equation}

Since $\hat{\theta}_n \xrightarrow{p} \theta$ (by assumption) and $R_n = o_p(|\hat{\theta}_n - \theta|)$, we have:
\begin{equation}
\sqrt{n}R_n = \sqrt{n} \cdot o_p(|\hat{\theta}_n - \theta|) = o_p(\sqrt{n}|\hat{\theta}_n - \theta|) = o_p(1)
\end{equation}

Therefore:
\begin{equation}
\sqrt{n}(g(\hat{\theta}_n) - g(\theta)) = g'(\theta)\sqrt{n}(\hat{\theta}_n - \theta) + o_p(1)
\end{equation}

By assumption, $\sqrt{n}(\hat{\theta}_n - \theta) \xrightarrow{d} \mathcal{N}(0, \sigma^2)$, and by continuous mapping theorem (multiplication by constant):
\begin{equation}
g'(\theta)\sqrt{n}(\hat{\theta}_n - \theta) \xrightarrow{d} \mathcal{N}(0, [g'(\theta)]^2\sigma^2)
\end{equation}

By Slutsky's theorem ($o_p(1) \to 0$), the result follows.
\end{proof}

\begin{example}[Application to Relative Lift]
For conversion rates $\hat{p}_1, \hat{p}_2$, the relative lift is:
\begin{equation}
g(p_1, p_2) = \frac{p_2 - p_1}{p_1} = \frac{p_2}{p_1} - 1
\end{equation}

By multivariate delta method:
\begin{equation}
\text{Var}\left(\frac{\hat{p}_2 - \hat{p}_1}{\hat{p}_1}\right) \approx \frac{1}{p_1^2}\text{Var}(\hat{p}_2) + \frac{p_2^2}{p_1^4}\text{Var}(\hat{p}_1)
\end{equation}
\end{example}

\section{Power Function Derivations}

The power function characterizes the probability of rejecting $H_0$ as a function of the true parameter value.

\subsection{Power Function for Two-Sample t-Test}

\begin{definition}[Power Function]
For testing $H_0: \mu_1 = \mu_2$ versus $H_1: \mu_1 \neq \mu_2$, the power function is:
\begin{equation}
\beta(\delta) = P(\text{Reject } H_0 \mid \mu_2 - \mu_1 = \delta)
\end{equation}
where $\delta$ is the true mean difference.
\end{definition}

\begin{theorem}[Power Function for t-Test]
For the two-sample t-test with equal sample sizes $n$ per group and common variance $\sigma^2$:
\begin{equation}
\beta(\delta) = \Phi\left(\frac{|\delta|\sqrt{n}}{\sigma\sqrt{2}} - z_{\alpha/2}\right) + \Phi\left(-\frac{|\delta|\sqrt{n}}{\sigma\sqrt{2}} - z_{\alpha/2}\right)
\end{equation}
where $\Phi$ is the standard normal CDF.
\end{theorem}

\begin{proof}
The test statistic is:
\begin{equation}
T = \frac{\bar{X}_2 - \bar{X}_1}{\sqrt{2\sigma^2/n}}
\end{equation}

Reject $H_0$ if $|T| > z_{\alpha/2}$, i.e., if $T > z_{\alpha/2}$ or $T < -z_{\alpha/2}$.

Under the alternative $\mu_2 - \mu_1 = \delta$:
\begin{equation}
\mathbb{E}[T \mid \delta] = \frac{\delta}{\sqrt{2\sigma^2/n}} = \frac{\delta\sqrt{n}}{\sigma\sqrt{2}}
\end{equation}

and $\text{Var}(T) = 1$. Therefore:
\begin{equation}
T \mid \delta \sim \mathcal{N}\left(\frac{\delta\sqrt{n}}{\sigma\sqrt{2}}, 1\right)
\end{equation}

The power is:
\begin{align}
\beta(\delta) &= P(T > z_{\alpha/2} \mid \delta) + P(T < -z_{\alpha/2} \mid \delta) \\
&= P\left(\mathcal{N}\left(\frac{\delta\sqrt{n}}{\sigma\sqrt{2}}, 1\right) > z_{\alpha/2}\right) + P\left(\mathcal{N}\left(\frac{\delta\sqrt{n}}{\sigma\sqrt{2}}, 1\right) < -z_{\alpha/2}\right)
\end{align}

Standardizing:
\begin{align}
&= P\left(\mathcal{N}(0,1) > z_{\alpha/2} - \frac{\delta\sqrt{n}}{\sigma\sqrt{2}}\right) + P\left(\mathcal{N}(0,1) < -z_{\alpha/2} - \frac{\delta\sqrt{n}}{\sigma\sqrt{2}}\right) \\
&= \Phi\left(\frac{\delta\sqrt{n}}{\sigma\sqrt{2}} - z_{\alpha/2}\right) + \Phi\left(-\frac{\delta\sqrt{n}}{\sigma\sqrt{2}} - z_{\alpha/2}\right)
\end{align}

For $\delta$ large, the second term is negligible, so:
\begin{equation}
\beta(\delta) \approx \Phi\left(\frac{|\delta|\sqrt{n}}{\sigma\sqrt{2}} - z_{\alpha/2}\right)
\end{equation}
\end{proof}

\subsection{Power Function for Proportion Test}

\begin{theorem}[Power for Two-Proportion Test]
For testing $H_0: p_1 = p_2$ versus $H_1: p_1 \neq p_2$ with sample size $n$ per group:
\begin{equation}
\beta(p_1, p_2) \approx \Phi\left(\frac{(p_2-p_1)\sqrt{n}}{\sqrt{p_1(1-p_1) + p_2(1-p_2)}} - z_{\alpha/2}\right)
\end{equation}
for $p_2 > p_1$.
\end{theorem}

\begin{proof}
Under $H_1$, the test statistic:
\begin{equation}
Z = \frac{\hat{p}_2 - \hat{p}_1}{\sqrt{[p_1(1-p_1) + p_2(1-p_2)]/n}} \sim \mathcal{N}\left(\frac{(p_2-p_1)\sqrt{n}}{\sqrt{p_1(1-p_1) + p_2(1-p_2)}}, 1\right)
\end{equation}

Reject if $|Z| > z_{\alpha/2}$ (using pooled variance approximation under $H_0$). The power is:
\begin{align}
\beta &= P(Z > z_{\alpha/2} \mid H_1) \\
&= \Phi\left(\frac{(p_2-p_1)\sqrt{n}}{\sqrt{p_1(1-p_1) + p_2(1-p_2)}} - z_{\alpha/2}\right)
\end{align}
\end{proof}

\subsection{Non-Centrality Parameter}

\begin{definition}[Non-Centrality Parameter]
For the two-sample t-test, the non-centrality parameter is:
\begin{equation}
\lambda = \frac{\delta}{\sigma}\sqrt{\frac{n}{2}}
\end{equation}
\end{definition}

\begin{theorem}[Power in Terms of Non-Central t]
The exact power function using the non-central t-distribution with $\nu = 2(n-1)$ degrees of freedom and non-centrality parameter $\lambda$ is:
\begin{equation}
\beta(\delta) = P(t_{\nu, \lambda} > t_{\nu, \alpha/2}) + P(t_{\nu, \lambda} < -t_{\nu, \alpha/2})
\end{equation}
where $t_{\nu, \lambda}$ denotes a non-central t random variable.
\end{theorem}

\begin{remark}
For large $n$, the normal approximation we derived earlier is excellent. For small $n$, the non-central t distribution should be used.
\end{remark}

\section{Bayesian Posterior Calculations}

Bayesian inference provides an alternative framework for A/B testing, incorporating prior knowledge.

\subsection{Beta-Binomial Model}

\begin{theorem}[Beta-Binomial Conjugacy]
Let $p \sim \text{Beta}(\alpha, \beta)$ be the prior distribution for a proportion. If we observe $X \sim \text{Binomial}(n, p)$ successes in $n$ trials, the posterior distribution is:
\begin{equation}
p \mid X \sim \text{Beta}(\alpha + X, \beta + n - X)
\end{equation}
\end{theorem}

\begin{proof}
By Bayes' theorem:
\begin{equation}
f(p \mid X) = \frac{f(X \mid p) f(p)}{f(X)} \propto f(X \mid p) f(p)
\end{equation}

The likelihood function is:
\begin{equation}
f(X \mid p) = \binom{n}{X} p^X(1-p)^{n-X} \propto p^X(1-p)^{n-X}
\end{equation}

The prior density is:
\begin{equation}
f(p) = \frac{\Gamma(\alpha + \beta)}{\Gamma(\alpha)\Gamma(\beta)} p^{\alpha-1}(1-p)^{\beta-1} \propto p^{\alpha-1}(1-p)^{\beta-1}
\end{equation}

The posterior is proportional to:
\begin{align}
f(p \mid X) &\propto p^X(1-p)^{n-X} \cdot p^{\alpha-1}(1-p)^{\beta-1} \\
&= p^{X+\alpha-1}(1-p)^{n-X+\beta-1}
\end{align}

This is the kernel of a $\text{Beta}(\alpha + X, \beta + n - X)$ distribution. The normalizing constant makes it a proper density:
\begin{equation}
f(p \mid X) = \frac{\Gamma(\alpha + \beta + n)}{\Gamma(\alpha + X)\Gamma(\beta + n - X)} p^{\alpha+X-1}(1-p)^{\beta+n-X-1}
\end{equation}
\end{proof}

\subsection{Posterior Predictive Distribution}

\begin{theorem}[Beta-Binomial Predictive]
Given posterior $p \mid X \sim \text{Beta}(\alpha + X, \beta + n - X)$, the predictive probability of success in the next trial is:
\begin{equation}
P(\tilde{X} = 1 \mid X) = \frac{\alpha + X}{\alpha + \beta + n}
\end{equation}
\end{theorem}

\begin{proof}
The posterior predictive distribution is:
\begin{align}
P(\tilde{X} = 1 \mid X) &= \int_0^1 P(\tilde{X} = 1 \mid p) f(p \mid X) \, dp \\
&= \int_0^1 p \cdot f(p \mid X) \, dp \\
&= \mathbb{E}[p \mid X]
\end{align}

For $p \sim \text{Beta}(\alpha', \beta')$, the mean is $\mathbb{E}[p] = \frac{\alpha'}{\alpha' + \beta'}$. Therefore:
\begin{equation}
P(\tilde{X} = 1 \mid X) = \frac{\alpha + X}{\alpha + X + \beta + n - X} = \frac{\alpha + X}{\alpha + \beta + n}
\end{equation}
\end{proof}

\subsection{Probability of Superiority}

\begin{theorem}[Closed Form for Beta Distributions]
If $p_A \sim \text{Beta}(\alpha_A, \beta_A)$ and $p_B \sim \text{Beta}(\alpha_B, \beta_B)$ independently, then:
\begin{equation}
P(p_B > p_A) = \sum_{i=0}^{\alpha_B-1} \frac{B(\alpha_A + i, \beta_A + \beta_B)}{(\beta_B + i) \cdot B(1+i, \beta_B) \cdot B(\alpha_A, \beta_A)}
\end{equation}
where $B(\cdot, \cdot)$ is the beta function.
\end{theorem}

\begin{proof}
By definition:
\begin{align}
P(p_B > p_A) &= \int_0^1 \int_0^{p_B} f_A(p_A) f_B(p_B) \, dp_A \, dp_B \\
&= \int_0^1 f_B(p_B) F_A(p_B) \, dp_B
\end{align}

where $F_A$ is the CDF of Beta$(\alpha_A, \beta_A)$.

The CDF of a Beta distribution can be expressed in terms of the incomplete beta function:
\begin{equation}
F_A(p) = I_p(\alpha_A, \beta_A) = \frac{B_p(\alpha_A, \beta_A)}{B(\alpha_A, \beta_A)}
\end{equation}

Substituting and using properties of beta integrals leads to the series representation above.

For computational purposes, when $\alpha_A, \alpha_B$ are integers (as in A/B testing with count data), this series is finite and easily computed.
\end{proof}

\subsection{Normal-Normal Conjugacy}

\begin{theorem}[Normal Prior, Normal Likelihood]
Suppose $\mu \sim \mathcal{N}(\mu_0, \tau^2)$ is the prior for a mean, and we observe data $\bar{X} \mid \mu \sim \mathcal{N}(\mu, \sigma^2/n)$ with known variance $\sigma^2$. Then the posterior is:
\begin{equation}
\mu \mid \bar{X} \sim \mathcal{N}\left(\frac{\frac{n}{\sigma^2}\bar{X} + \frac{1}{\tau^2}\mu_0}{\frac{n}{\sigma^2} + \frac{1}{\tau^2}}, \left(\frac{n}{\sigma^2} + \frac{1}{\tau^2}\right)^{-1}\right)
\end{equation}
\end{theorem}

\begin{proof}
The posterior is proportional to likelihood times prior:
\begin{align}
f(\mu \mid \bar{X}) &\propto \exp\left(-\frac{n(\bar{X} - \mu)^2}{2\sigma^2}\right) \cdot \exp\left(-\frac{(\mu - \mu_0)^2}{2\tau^2}\right) \\
&\propto \exp\left(-\frac{1}{2}\left[\frac{n(\bar{X} - \mu)^2}{\sigma^2} + \frac{(\mu - \mu_0)^2}{\tau^2}\right]\right)
\end{align}

Expand the exponent:
\begin{align}
&\frac{n(\bar{X} - \mu)^2}{\sigma^2} + \frac{(\mu - \mu_0)^2}{\tau^2} \\
&= \frac{n}{\sigma^2}(\bar{X}^2 - 2\bar{X}\mu + \mu^2) + \frac{1}{\tau^2}(\mu^2 - 2\mu_0\mu + \mu_0^2) \\
&= \left(\frac{n}{\sigma^2} + \frac{1}{\tau^2}\right)\mu^2 - 2\left(\frac{n\bar{X}}{\sigma^2} + \frac{\mu_0}{\tau^2}\right)\mu + \text{const}
\end{align}

Let $\lambda = \frac{n}{\sigma^2} + \frac{1}{\tau^2}$ (posterior precision). Complete the square:
\begin{align}
&= \lambda\left[\mu^2 - 2\mu \cdot \frac{1}{\lambda}\left(\frac{n\bar{X}}{\sigma^2} + \frac{\mu_0}{\tau^2}\right)\right] + \text{const} \\
&= \lambda\left(\mu - \frac{1}{\lambda}\left(\frac{n\bar{X}}{\sigma^2} + \frac{\mu_0}{\tau^2}\right)\right)^2 + \text{const}
\end{align}

This is the kernel of $\mathcal{N}(\mu_n, \sigma_n^2)$ where:
\begin{align}
\mu_n &= \frac{1}{\lambda}\left(\frac{n\bar{X}}{\sigma^2} + \frac{\mu_0}{\tau^2}\right) = \frac{\frac{n}{\sigma^2}\bar{X} + \frac{1}{\tau^2}\mu_0}{\frac{n}{\sigma^2} + \frac{1}{\tau^2}} \\
\sigma_n^2 &= \frac{1}{\lambda} = \left(\frac{n}{\sigma^2} + \frac{1}{\tau^2}\right)^{-1}
\end{align}
\end{proof}

\begin{remark}[Interpretation]
The posterior mean is a precision-weighted average of the prior mean $\mu_0$ and the data mean $\bar{X}$. As $n \to \infty$, the posterior converges to $\mathcal{N}(\bar{X}, \sigma^2/n)$, i.e., the data overwhelms the prior.
\end{remark}

\section{Sequential Testing Boundaries}

Sequential testing allows early stopping based on accumulating evidence, while controlling error rates.

\subsection{Sequential Probability Ratio Test (SPRT)}

\begin{theorem}[Wald's SPRT]
For testing $H_0: \theta = \theta_0$ versus $H_1: \theta = \theta_1$, the sequential probability ratio test stops at time $n$ when:
\begin{equation}
\Lambda_n = \frac{L(\theta_1; X_1, \ldots, X_n)}{L(\theta_0; X_1, \ldots, X_n)} \notin (B, A)
\end{equation}
and rejects $H_0$ if $\Lambda_n \geq A$, accepts $H_0$ if $\Lambda_n \leq B$, where:
\begin{equation}
A = \frac{1-\beta}{\alpha}, \quad B = \frac{\beta}{1-\alpha}
\end{equation}
for desired error rates $\alpha$ (Type I) and $\beta$ (Type II).
\end{theorem}

\begin{proof}[Proof Sketch]
At stopping time $\tau$, if we reject $H_0$:
\begin{equation}
\frac{P(X_1, \ldots, X_\tau \mid \theta_1)}{P(X_1, \ldots, X_\tau \mid \theta_0)} \geq A
\end{equation}

This implies:
\begin{equation}
P(X_1, \ldots, X_\tau \mid \theta_1) \geq A \cdot P(X_1, \ldots, X_\tau \mid \theta_0)
\end{equation}

Summing over all sequences leading to rejection:
\begin{equation}
P(\text{Reject} \mid \theta_1) \geq A \cdot P(\text{Reject} \mid \theta_0)
\end{equation}

Therefore:
\begin{equation}
P(\text{Reject} \mid \theta_0) \leq \frac{P(\text{Reject} \mid \theta_1)}{A} = \frac{1-\beta}{A}
\end{equation}

Wald showed that choosing $A = (1-\beta)/\alpha$ achieves approximately $P(\text{Reject} \mid \theta_0) \approx \alpha$.
\end{proof}

\subsection{Group Sequential Boundaries}

For group sequential designs with $K$ planned looks, we derive boundaries that maintain overall Type I error rate $\alpha$.

\begin{theorem}[O'Brien-Fleming Boundaries]
For $K$ interim analyses at information fractions $t_1, \ldots, t_K$ (where $t_K = 1$), the O'Brien-Fleming boundaries are:
\begin{equation}
c_k = \frac{c}{\sqrt{t_k}}, \quad k = 1, \ldots, K
\end{equation}
where $c$ is chosen such that:
\begin{equation}
P\left(\max_{1 \leq k \leq K} \frac{Z_k}{\sqrt{t_k}} > c \mid H_0\right) = \alpha
\end{equation}
and $Z_k$ is the standardized test statistic at analysis $k$.
\end{theorem}

\begin{proof}[Derivation]
Let $Z_1, \ldots, Z_K$ be the sequence of standardized test statistics with information times $t_1 < \cdots < t_K = 1$.

Under $H_0$, by the CLT and properties of independent increments:
\begin{equation}
(Z_1, \ldots, Z_K) \text{ follows a multivariate normal with } \text{Cov}(Z_i, Z_j) = \sqrt{t_i/t_j} \text{ for } i \leq j
\end{equation}

The O'Brien-Fleming approach standardizes by $\sqrt{t_k}$ at each look:
\begin{equation}
W_k = \frac{Z_k}{\sqrt{t_k}}
\end{equation}

We reject $H_0$ at the first $k$ where $|W_k| > c$.

The Type I error is:
\begin{equation}
\alpha = P\left(\bigcup_{k=1}^K \{|W_k| > c\}\right) = P\left(\max_{1 \leq k \leq K} |W_k| > c\right)
\end{equation}

For equal spacing and $K$ looks, approximate values of $c$ are:
\begin{itemize}
    \item $K=2$: $c \approx 2.24$ (vs. 1.96 for fixed design)
    \item $K=5$: $c \approx 2.04$
    \item $K \to \infty$: $c \to 1.96$
\end{itemize}

These are computed numerically to achieve exact $\alpha$ (typically $\alpha = 0.05$).
\end{proof}

\subsection{Pocock Boundaries}

\begin{theorem}[Pocock Boundaries]
The Pocock boundaries use a constant critical value at each analysis:
\begin{equation}
c_k = c, \quad k = 1, \ldots, K
\end{equation}
where $c$ satisfies:
\begin{equation}
P\left(\max_{1 \leq k \leq K} Z_k > c \mid H_0\right) = \alpha
\end{equation}
\end{theorem}

\begin{remark}[Comparison]
\begin{itemize}
    \item O'Brien-Fleming: Early boundaries are very high (conservative early, liberal late). Preserves power if test goes to completion.
    \item Pocock: Constant boundaries (liberal early, conservative late). Higher early stopping probability.
\end{itemize}
\end{remark}

\subsection{Alpha Spending Functions}

\begin{definition}[Alpha Spending Function]
An alpha spending function $\alpha(t)$ for $t \in [0, 1]$ satisfies:
\begin{enumerate}
    \item $\alpha(0) = 0$
    \item $\alpha(1) = \alpha$ (overall Type I error)
    \item $\alpha(t)$ is non-decreasing
\end{enumerate}

At analysis $k$ with information time $t_k$, the incremental alpha is:
\begin{equation}
\alpha_k = \alpha(t_k) - \alpha(t_{k-1})
\end{equation}
\end{definition}

\begin{example}[Common Spending Functions]
\textbf{O'Brien-Fleming type:}
\begin{equation}
\alpha(t) = 2\left[1 - \Phi\left(\frac{z_{\alpha/2}}{\sqrt{t}}\right)\right]
\end{equation}

\textbf{Pocock type:}
\begin{equation}
\alpha(t) = \alpha \log(1 + (e-1)t)
\end{equation}

\textbf{Power family:}
\begin{equation}
\alpha(t) = \alpha t^\rho, \quad \rho > 0
\end{equation}
\end{example}

\subsection{Information-Based Monitoring}

\begin{theorem}[Conditional Error Principle]
At interim analysis $k$ with observed test statistic $Z_k$ and information fraction $t_k$, the boundary $c_k$ is determined by:
\begin{equation}
P\left(\bigcup_{j=k}^K \{Z_j > c_j\} \mid Z_k \right) = \frac{\alpha(t_k) - \alpha(t_{k-1})}{1 - \alpha(t_{k-1})}
\end{equation}
given that $H_0$ was not rejected at previous looks.
\end{theorem}

\begin{proof}[Derivation Sketch]
This follows from the law of total probability and the requirement that:
\begin{equation}
P(\text{Reject at or after look } k \mid H_0) = \alpha(t_k)
\end{equation}

Conditioning on not having rejected before look $k$:
\begin{align}
\alpha(t_k) &= P(\text{Reject at look } k \mid H_0) + P(\text{Reject after look } k \mid H_0) \\
&= P(\text{Reject at look } k \mid H_0, \text{Not rejected before}) \cdot P(\text{Not rejected before}) \\
&\quad + P(\text{Reject after look } k \mid H_0, \text{Not rejected at } k) \cdot P(\text{Not rejected at or before } k)
\end{align}

Working through the conditional probabilities yields the stated formula for determining $c_k$.
\end{proof}

\section{Additional Results}

\subsection{Variance of Sample Proportion}

\begin{theorem}
For $\hat{p} = X/n$ where $X \sim \text{Binomial}(n, p)$:
\begin{equation}
\text{Var}(\hat{p}) = \frac{p(1-p)}{n}
\end{equation}
\end{theorem}

\begin{proof}
Since $X \sim \text{Binomial}(n, p)$, we have $\text{Var}(X) = np(1-p)$. Therefore:
\begin{equation}
\text{Var}(\hat{p}) = \text{Var}\left(\frac{X}{n}\right) = \frac{1}{n^2}\text{Var}(X) = \frac{1}{n^2} \cdot np(1-p) = \frac{p(1-p)}{n}
\end{equation}
\end{proof}

\subsection{CUPED Variance Reduction}

\begin{theorem}[CUPED Estimator]
For outcome $Y$ and covariate $X$, the CUPED-adjusted estimator:
\begin{equation}
Y^* = Y - \theta(X - \mathbb{E}[X])
\end{equation}
with optimal $\theta = \text{Cov}(Y, X)/\text{Var}(X)$ achieves variance:
\begin{equation}
\text{Var}(Y^*) = \text{Var}(Y)(1 - \rho^2)
\end{equation}
where $\rho = \text{Cor}(Y, X)$ is the correlation coefficient.
\end{theorem}

\begin{proof}
Expanding the variance:
\begin{align}
\text{Var}(Y^*) &= \text{Var}(Y - \theta(X - \mathbb{E}[X])) \\
&= \text{Var}(Y) + \theta^2\text{Var}(X) - 2\theta\text{Cov}(Y, X)
\end{align}

Substituting $\theta = \text{Cov}(Y, X)/\text{Var}(X)$:
\begin{align}
\text{Var}(Y^*) &= \text{Var}(Y) + \frac{\text{Cov}(Y,X)^2}{\text{Var}(X)^2}\text{Var}(X) - 2\frac{\text{Cov}(Y,X)}{\text{Var}(X)}\text{Cov}(Y, X) \\
&= \text{Var}(Y) + \frac{\text{Cov}(Y,X)^2}{\text{Var}(X)} - \frac{2\text{Cov}(Y,X)^2}{\text{Var}(X)} \\
&= \text{Var}(Y) - \frac{\text{Cov}(Y,X)^2}{\text{Var}(X)}
\end{align}

Since $\rho = \text{Cov}(Y,X)/(\sigma_Y\sigma_X)$, we have $\text{Cov}(Y,X)^2 = \rho^2\text{Var}(Y)\text{Var}(X)$:
\begin{align}
\text{Var}(Y^*) &= \text{Var}(Y) - \frac{\rho^2\text{Var}(Y)\text{Var}(X)}{\text{Var}(X)} \\
&= \text{Var}(Y)(1 - \rho^2)
\end{align}
\end{proof}

\subsection{Bonferroni and Šidák Corrections}

\begin{theorem}[Bonferroni Inequality]
For $m$ hypothesis tests, each conducted at level $\alpha/m$, the family-wise error rate satisfies:
\begin{equation}
\text{FWER} \leq \alpha
\end{equation}
\end{theorem}

\begin{proof}
Let $R_i$ be the event that test $i$ rejects a true null hypothesis. By Boole's inequality (union bound):
\begin{align}
\text{FWER} &= P\left(\bigcup_{i=1}^m R_i\right) \\
&\leq \sum_{i=1}^m P(R_i) \\
&= \sum_{i=1}^m \frac{\alpha}{m} = \alpha
\end{align}
\end{proof}

\begin{theorem}[Šidák Correction]
If the $m$ tests are independent, testing each at level $\alpha_s = 1 - (1-\alpha)^{1/m}$ gives exact FWER $= \alpha$.
\end{theorem}

\begin{proof}
Under independence:
\begin{align}
\text{FWER} &= P(\text{Reject at least one true } H_0) \\
&= 1 - P(\text{Accept all true } H_0) \\
&= 1 - \prod_{i=1}^m P(\text{Accept } H_{0i}) \\
&= 1 - \prod_{i=1}^m (1 - \alpha_s) \\
&= 1 - (1 - \alpha_s)^m
\end{align}

Setting this equal to $\alpha$:
\begin{align}
1 - (1-\alpha_s)^m &= \alpha \\
(1-\alpha_s)^m &= 1 - \alpha \\
\alpha_s &= 1 - (1-\alpha)^{1/m}
\end{align}
\end{proof}

\begin{remark}
The Šidák correction is less conservative than Bonferroni when tests are independent, since $(1-\alpha)^{1/m} > 1 - \alpha/m$ for $m > 1$.
\end{remark}

\section{Summary}

This appendix has provided rigorous mathematical foundations for:

\begin{itemize}
    \item \textbf{Sample size formulas:} Derived for t-tests and proportion tests, including unequal allocation
    \item \textbf{Central Limit Theorem:} Proof via characteristic functions with Berry-Esseen bound
    \item \textbf{Power functions:} Explicit derivations showing power as function of effect size
    \item \textbf{Bayesian methods:} Beta-binomial and normal-normal conjugacy with posterior calculations
    \item \textbf{Sequential testing:} SPRT, O'Brien-Fleming, and Pocock boundaries with alpha spending functions
\end{itemize}

These mathematical results underpin the practical A/B testing methods presented throughout this handbook.
